% Generated draft: requires mathematical and source review.
\documentclass[11pt]{article}
\usepackage[margin=25mm]{geometry}
\usepackage{fontspec}
\setmainfont{Latin Modern Roman}
\usepackage{amsmath,amssymb,mathtools,unicode-math}
\setmathfont{Latin Modern Math}
\usepackage{enumitem,needspace,xurl,hyperref}
\setlength{\parindent}{0pt}
\setlength{\parskip}{5pt}
\setlength{\emergencystretch}{3em}
\newcommand{\N}{\mathbb{N}}
\newcommand{\Z}{\mathbb{Z}}
\newcommand{\Q}{\mathbb{Q}}
\newcommand{\R}{\mathbb{R}}
\newcommand{\C}{\mathbb{C}}
\newcommand{\F}{\mathbb{F}}
\newcommand{\A}{\mathbb{A}}
\newcommand{\PP}{\mathbb{P}}
\newcommand{\E}{\mathbb{E}}
\newcommand{\eps}{\varepsilon}
\newcommand{\dd}{\,\mathrm{d}}
\begin{document}
\section*{UnsolvedMath ID: 701569}
\textbf{Source problem code:} catalogue entry
\subsection{Definitions and mathematical statement}
Use the standard convention that a locally compact group is a Hausdorff topological group whose underlying space is locally compact; the supplied statement does not explicitly specify a separation convention. Let \(e\) denote its identity. Such a group \(G\) is compactly generated if there is a compact subset \(K\subseteq G\) with \(\langle K\rangle=G\), and is topologically simple if its only closed normal subgroups are \(\{e\}\) and \(G\).

For a left Haar measure \(\mu\) on \(G\), amenability means that there exists a positive linear functional \(m:L^\infty(G,\mu;\R)\to\R\) such that \(m(1)=1\) and
\[
 m(L_g f)=m(f)\qquad\text{for all }g\in G\text{ and }f\in L^\infty(G,\mu;\R),
 \qquad (L_g f)(x)=f(g^{-1}x).
\]
Here \(L^\infty(G,\mu;\R)\) consists of essentially bounded measurable real functions modulo equality almost everywhere, and positivity means \(m(f)\geq0\) whenever \(f\geq0\) almost everywhere.

\textbf{Question.} Does there exist a locally compact group \(G\) that is compactly generated, topologically simple, amenable, non-discrete, and non-compact? Non-discreteness means that \(\{e\}\) is not open in \(G\); non-compactness refers to the underlying topological space.

The source states that a positive answer implies a negative answer to Kourovka Notebook Problem 19.107. The statement of Problem 19.107 is not supplied, so this implication is retained as a source qualification without a formal restatement or independent verification.
\subsection{Short English statement}
Does there exist a non-discrete, non-compact, compactly generated locally compact group that is both amenable and topologically simple? This is Problem 19.21, proposed by P.-E. Caprace in 2018. The local record labels it open; that label is supplied curation, not an independently verified assessment.
\subsection{Sources}
\begin{itemize}
\item Local UnsolvedMath registry, permanent ID \(701569\), exact source problem code \texttt{catalogue entry}, titled \emph{Kourovka Notebook Problem 19.21}. The supplied record labels the problem open.
\item \emph{The Kourovka Notebook (Unsolved Problems in Group Theory)}, Problem 19.21, proposed by P.-E. Caprace; the supplied background identifies the 19th issue (2018) and the 21st issue edition (2026), arXiv:1401.0300v48. The collection description names E. I. Khukhro and V. D. Mazurov as editors.
\item The supplied statement asserts that a positive answer implies a negative answer to Problem 19.107. That problem's statement and a justification of the implication are not included in the record. No independent source consultation is claimed.
\end{itemize}
\end{document}
